\documentclass[10pt,landscape]{article}

\usepackage[margin=0.32in]{geometry}
\usepackage[T1]{fontenc}
\usepackage[utf8]{inputenc}
\usepackage{lmodern}
\usepackage{textcomp}
\usepackage{microtype}
\usepackage[table,dvipsnames]{xcolor}
\usepackage{array}
\usepackage{booktabs}
\usepackage{enumitem}
\usepackage{hyperref}
\usepackage{listings}
\usepackage{multicol}
\usepackage{tabularx}
\usepackage{titlesec}

\pagestyle{empty}
\setlength{\parindent}{0pt}
\setlength{\parskip}{1.5pt}
\setlength{\columnsep}{13pt}
\setlength{\columnseprule}{0.2pt}
\setcounter{secnumdepth}{0}

\definecolor{ToshBlue}{HTML}{0F4C81}
\definecolor{ToshInk}{HTML}{1F2937}
\definecolor{ToshMuted}{HTML}{5B6470}
\definecolor{ToshPanel}{HTML}{F4F7FB}
\definecolor{ToshRule}{HTML}{CBD5E1}
\definecolor{ToshCodeBg}{HTML}{F8FAFC}

\color{ToshInk}

\titleformat{\section}
  {\large\bfseries\color{ToshBlue}}
  {}
  {0pt}
  {}
\titlespacing*{\section}{0pt}{0.4em}{0.12em}

\setlist[itemize]{leftmargin=1.05em,itemsep=0.5pt,topsep=1.5pt,parsep=0pt,partopsep=0pt}

\lstdefinestyle{tosh}{
  basicstyle=\ttfamily\scriptsize,
  backgroundcolor=\color{ToshCodeBg},
  frame=single,
  rulecolor=\color{ToshRule},
  framerule=0.4pt,
  xleftmargin=0pt,
  xrightmargin=0pt,
  aboveskip=0.18em,
  belowskip=0.24em,
  columns=fullflexible,
  keepspaces=true,
  showstringspaces=false,
  breaklines=true,
  literate={°}{{\textdegree}}1
           {μ}{{\ensuremath{\mu}}}1
           {Ω}{{\ensuremath{\Omega}}}1
}
\lstset{style=tosh}

\newcommand{\sheettitle}[3]{
  {\LARGE\bfseries\color{ToshBlue} ToSh Cheat Sheet: #1\par}
  \vspace{0.08em}
  {\small\color{ToshMuted} #2 \hfill #3\par}
  \vspace{0.18em}
  {\color{ToshRule}\rule{\linewidth}{0.8pt}}
  \vspace{0.12em}
}

\newcommand{\note}[1]{
  \vspace{0.12em}
  \fcolorbox{ToshRule}{ToshPanel}{%
    \parbox{\dimexpr\linewidth-2\fboxsep-2\fboxrule\relax}{\footnotesize #1}%
  }
  \vspace{0.12em}
}

\newcommand{\tighttable}[2]{
  {\footnotesize
  \renewcommand{\arraystretch}{1.08}
  \begin{tabularx}{\linewidth}{@{}#1@{}}
  #2
  \end{tabularx}}
}


\begin{document}
\small
\sheettitle{Filesystem And File I/O}{Navigation, discovery, mutation, text/binary commands, and managed handles}{command reference + tests}

\begin{multicols*}{3}

\section{Navigation}
\tighttable{>{\ttfamily\raggedright\arraybackslash}p{0.28\linewidth}X}{
\toprule
Command & Use \\
\midrule
\texttt{pwd} & current directory as an object \\
\texttt{cd path} & change directory \\
\texttt{dirs} & inspect the directory stack \\
\texttt{back} & move backward in the stack \\
\texttt{forward} & move forward in the stack \\
\bottomrule
}

\begin{lstlisting}
pwd
cd src
dirs
back
forward
\end{lstlisting}

\section{Discovery}
\begin{lstlisting}
ls -la
ls -R --group-directories-first
stat README.md
tree -L 2
\end{lstlisting}

\begin{lstlisting}
find . -name *.tosh
find . -type dir -maxdepth 2
glob **/*.tosh
realpath ./README.md
readlink ./some-link
\end{lstlisting}

\tighttable{>{\ttfamily\raggedright\arraybackslash}p{0.31\linewidth}X}{
\toprule
Predicate & Meaning \\
\midrule
\texttt{exists} & path exists \\
\texttt{is-file} & regular file check \\
\texttt{is-dir} & directory check \\
\texttt{is-link} & symlink check \\
\texttt{basename} & file name only \\
\texttt{dirname} & parent path only \\
\bottomrule
}

\section{Typed File Rows}
\begin{lstlisting}
ls -la | where _.Type == file | get { Name, Size, Modified }
find . -name *.cs | first 10 | get FullName
stat README.md | inspect
\end{lstlisting}

\note{Filesystem commands emit typed entries, so fields like \texttt{Name}, \texttt{Type}, \texttt{Size}, and \texttt{Modified} can be sorted or filtered directly.}

\columnbreak

\section{Mutation}
\begin{lstlisting}
mkdir scratch
touch scratch/notes.txt
cp README.md scratch/readme-copy.md
mv scratch/readme-copy.md scratch/info.md
rm scratch/info.md
\end{lstlisting}

\begin{lstlisting}
chmod 755 ./script
ln -s target.txt link.txt
mkdir-temp
tempfile
\end{lstlisting}

\tighttable{>{\ttfamily\raggedright\arraybackslash}p{0.29\linewidth}X}{
\toprule
Command & Notes \\
\midrule
\texttt{mkdir} & supports parent creation flags \\
\texttt{touch} & create or update timestamps \\
\texttt{cp}, \texttt{mv}, \texttt{rm} & typed file operations \\
\texttt{chmod}, \texttt{chown}, \texttt{ln} & Unix-style metadata/link tools \\
\bottomrule
}

\section{Read And Write Files}
\begin{lstlisting}
cat README.md
read-file README.md
read-lines README.md | first 5
read-bytes icon.png
\end{lstlisting}

\begin{lstlisting}
write-file scratch/notes.txt "hello"
append-file scratch/notes.txt "world"
write-bytes scratch/blob.bin [1, 2, 3]
\end{lstlisting}

\tighttable{>{\ttfamily\raggedright\arraybackslash}p{0.30\linewidth}X}{
\toprule
Command & Shape \\
\midrule
\texttt{read-file} & one file -> one string \\
\texttt{read-lines} & lines as separate values \\
\texttt{read-bytes} & byte arrays \\
\texttt{lines} / \texttt{join-lines} & text stream conversion \\
\texttt{as-file} & materialize pipeline output as a temp file \\
\bottomrule
}

\section{Managed Handles}
\begin{lstlisting}
var reader = (open-file --read scratch/notes.txt)
read-line-from $reader
read-to-end $reader
close $reader
\end{lstlisting}

\begin{lstlisting}
var writer = (open-file --write scratch/log.txt)
write-line-to $writer "ready"
flush $writer
close $writer
\end{lstlisting}

\columnbreak

\section{Handle Toolbox}
\tighttable{>{\ttfamily\raggedright\arraybackslash}p{0.32\linewidth}X}{
\toprule
Command & Use \\
\midrule
\texttt{open-file} & create text or binary file handles \\
\texttt{read-from} & read a chunk from a handle \\
\texttt{read-line-from} & read the next text line \\
\texttt{read-to-end} & consume the rest of the handle \\
\texttt{write-to} / \texttt{write-line-to} & write bytes or text \\
\texttt{seek}, \texttt{position}, \texttt{length} & stream positioning \\
\texttt{copy-to} & copy one handle into another \\
\texttt{flush}, \texttt{close} & finish safely \\
\bottomrule
}

\section{Useful Recipes}
\begin{lstlisting}
# biggest files here
ls -la | where _.Type == file | sort Size | reverse | first 10

# recent logs
find . -name *.log | where _.Modified > ((date now) - 1d)
\end{lstlisting}

\begin{lstlisting}
# line-oriented pipeline
read-lines README.md | grep "ToSh" | first 5

# stage pipeline output into a file for externals
var names = (ls | get Name | sort | as-file text)
/bin/cat $names
\end{lstlisting}

\section{External-System Wrappers}
\begin{lstlisting}
df
du .
lsblk -l
findmnt
systemctl --type service | first 5
ip addr
\end{lstlisting}

\begin{itemize}
\item Many system commands are wrapped and return structured rows instead of raw text.
\item That means \texttt{sort}, \texttt{where}, and \texttt{summarize} work on their fields directly.
\end{itemize}

\note{\textbf{Good habit:} prefer ToSh wrappers like \texttt{ls}, \texttt{find}, \texttt{df}, \texttt{ip}, and \texttt{systemctl} when you want typed data; drop to native tools when you need exact external behavior.}

\end{multicols*}
\newpage

\sheettitle{Filesystem And File I/O}{Typed wrappers, stream positioning, and path-heavy recipes}{page 2 / front-back reference}

\begin{multicols*}{3}

\section{Disk And Mount Views}
\begin{lstlisting}
df --total
df -t ext4 --show FileSystem,Type,UsePercent,MountedOn
df . | get { FileSystem, MountedOn }
\end{lstlisting}

\begin{lstlisting}
du -s .
du -a -c --time
du -x ./projects | get { Name, Size, Modified }
\end{lstlisting}

\note{\texttt{df} and \texttt{du} return typed size fields, so sorting and summarizing disk usage stays object-oriented instead of text-oriented.}

\section{Structured System Wrappers}
\begin{lstlisting}
findmnt -l | where _.Target.StartsWith("/run")
lsblk -l -f | where _.FsType == "ntfs"
lsfd --summary=only
\end{lstlisting}

\begin{lstlisting}
ip addr | each { _.Addresses } | flatten | get Address
systemctl status sshd.service | get { Id, ActiveState, MainPID, RecentLogCount }
journalctl -u sshd.service | first 10
\end{lstlisting}

\columnbreak

\section{Path Utilities}
\begin{lstlisting}
glob -a **/*.cs
realpath ./relative/path
readlink -f ./chain
basename /home/user/file.txt
dirname /home/user/file.txt
\end{lstlisting}

\tighttable{>{\ttfamily\raggedright\arraybackslash}p{0.30\linewidth}X}{
\toprule
Tool & When to use it \\
\midrule
\texttt{glob} & shell-style wildcard expansion into typed paths \\
\texttt{realpath} & canonicalize a relative path \\
\texttt{readlink} & inspect a symlink target or whole chain \\
\texttt{basename} / \texttt{dirname} & split path text safely \\
\bottomrule
}

\section{Handle Positioning}
\begin{lstlisting}
var h = (open-file --binary ./data.bin)
length $h
position $h
seek $h 128 current
read-from $h 32
\end{lstlisting}

\begin{lstlisting}
var src = (open-file --read ./a.txt)
var dst = (open-file --write ./b.txt)
copy-to $src $dst
close $src $dst
\end{lstlisting}

\note{\texttt{seek} returns the handle again, so you can keep piping through \texttt{read-from}, \texttt{read-to-end}, or \texttt{copy-to}.}

\section{Copy, Link, And Mode Tricks}
\begin{lstlisting}
cp -u notes.txt backup.txt
cp -s file.txt symlink.txt
cp -l file.txt hardlink.txt
chmod +x script.sh
chmod u+rw,go-w file.txt
\end{lstlisting}

\columnbreak

\section{Write Strategies}
\begin{lstlisting}
echo alpha beta | write-file ./notes.txt
echo tail line | append-file ./notes.txt
read-bytes ./source.bin | write-bytes ./copy.bin
\end{lstlisting}

\begin{lstlisting}
var log = (open-file --append ./run.log)
write-line-to $log "ready"
flush $log
close $log
\end{lstlisting}

\section{Practical File Recipes}
\begin{lstlisting}
find . -name *.log | where _.Size >= 10mb | get { Name, Size }
glob **/*.md | where _.Modified > ((date now) - 7d)
ls -la | where _.Type == file | sort Modified | reverse | first 20
\end{lstlisting}

\begin{lstlisting}
var tmp = (mkdir-temp report)
var out = (tempfile export .csv)
$rows | to csv | write-file $out
realpath $out
\end{lstlisting}

\section{Mental Model}
\begin{itemize}
\item Path-first wrappers like \texttt{ls}, \texttt{stat}, \texttt{find}, \texttt{df}, and \texttt{lsblk} try hard to keep you in typed objects.
\item Managed handles are explicit on purpose: you can inspect them, pass them around, and close them deterministically.
\item Reach for external programs when you need exact native semantics; reach for ToSh wrappers when you want composable shell objects.
\end{itemize}

\end{multicols*}
\end{document}
